% !TITLE 天净沙·秋
% !AUTHOR 白朴
% !DYNASTY 元
% !GENRE 散曲
% !ORDER 3

\begin{poem}
孤村落日残霞\textsuperscript{1}，轻烟老树寒鸦\textsuperscript{2}，一点飞鸿影下\textsuperscript{3}。
青山绿水，白草红叶黄花\textsuperscript{4}。
\end{poem}

\begin{annotations}
\notetext{1}{残霞：即将消散的晚霞。与“落日”照应，点出黄昏时分。}
\notetext{2}{寒鸦：秋寒中的乌鸦。白朴的“寒鸦”与马致远的“昏鸦”着眼不同——一个强调寒冷，一个强调昏暗。}
\notetext{3}{飞鸿：大雁。鸿雁南飞是秋天的典型意象。“影下”写雁群飞过，影子掠过地面。}
\notetext{4}{白草红叶黄花：白色的草（经霜后枯萎变白）、红色的枫叶、黄色的菊花——三种颜色概括了秋天的多样色彩。}
\end{annotations}

\begin{appreciation}
《天净沙·秋》是白朴写的小令，和《天净沙·秋思》共用同一个曲牌，也是元代写秋天最有名的小令之一。

全曲只有两句，二十八个字。前半句：孤村落日残霞，轻烟老树寒鸦，一点飞鸿影下——一座孤零零的村子，落日，快要散尽的晚霞；轻烟，老树，秋天里冷得瑟缩的乌鸦；忽然一点影子掠过地面，是一行大雁飞过去了。注意这"一点飞鸿影下"：前面的景物全是静的，落日、残霞、老树、寒鸦，一样一样摆在那里不动；大雁一来，画面一下子活了。秋天最动人的时刻就在这种地方——万籁俱寂里突然动了一下，然后重新静下去。

后半句是全曲最妙的转折：青山绿水，白草红叶黄花。前面是黄昏的冷色调，这里忽然冒出五种颜色：青山，绿水，白色的草，红色的枫叶，黄色的菊花。寒鸦的"寒"和红叶的"红"，原来住在同一个秋天里。

细看这首和《天净沙·秋思》，像一对兄弟。同一个曲牌，同样是名词一件一件摆出来，同样停在黄昏。但马致远的主语是"断肠人"，白朴的主语是颜色本身。马致远的秋天是心情，白朴的秋天是风景；一个让你想家，一个让你抬头。

白朴是金朝大臣的儿子，金灭亡的时候他还是个孩子，一辈子没有做元朝的官，只在山水和曲子里过完一生。他笔下的秋天没有悲伤。也许不是他不懂，而是他选择只写好看的那一半——这选择本身，就像一点飞鸿影下，轻轻一下，却让人记得很久。

这一首，是哥哥希望你停下来看看的那一种。秋天放学路上，你抬头看一眼：天很高，柿子挂在枝头，银杏落了一地，跟曲子里的颜色一模一样。走路的时候别光顾着低头看手机，秋天不等人。
\end{appreciation}
